\documentclass[UTF8,fontset=none,11pt,a4paper]{ctexart}
\usepackage{ipara-handbook}
\handbooksetup{NET-B04 Flow × Congestion Control}{408 · Network Internal Bridge}{rwnd · cwnd · FlightSize · Feedback}
\begin{document}
\handbooktitle{Flow Control × Congestion Control}{两类反馈怎样共同限制发送}{408 · NET-B04}

\begin{corebox}{Mother Interface}
\[
\boxed{Usable=\max\{0,\min(rwnd,cwnd)-FlightSize\}}
\]
\code{rwnd} 是接收端容量声明，\code{cwnd} 是发送端对路径承载能力的控制状态；共同作用点是尚未确认的在途字节，但反馈源、保护对象与恢复事件不同。
\end{corebox}

\begin{methodbox}{Owners 与停止边界}
NET06 Owns receive buffer、advertised window 与 TCP 连接状态；NET07 Owns congestion evidence 和 \code{cwnd} 算法。本册只 Owns 两个回路的合成、辨识和同时变化时的推演。可靠字节区间交给 NET-B03，BDP 利用率交给 NET-B05。
\end{methodbox}

\section{两个回路必须各画一个 Owner}

接收端根据 buffer capacity、已缓存未被应用读取的数据计算可接收空间，并在 ACK 的 Window 字段中通告 \code{rwnd}。应用读取释放空间，数据到达消耗空间。路径不会替接收端生成这个值。

发送端根据 ACK 到达、loss/ECN/timeout 等路径证据维护 \code{cwnd}。接收端即使 buffer 很大，也不能授权发送端越过自身拥塞窗口。因而：
\[
\boxed{\text{Receiver loop protects endpoint}}
\qquad
\boxed{\text{Network loop protects path/queues}}.
\]

\section{合成公式中的三个量不能漏}

$\min(rwnd,cwnd)$ 是允许在途上限，不是“此刻还能新发多少”。若 $rwnd=12$ KB、$cwnd=8$ KB、已有 6 KB 未确认，则最多再发 2 KB，而不是 8 KB。ACK 推进会减小 FlightSize 并释放发送空间；窗口值变化会改变上限，二者是不同事件。

接收端通告 zero window 时，正常新数据停止，但 persist/zero-window probe 用于恢复丢失的 window update；这不是拥塞超时。反之，\code{cwnd} 因 loss 缩小时，不能据此推出 receiver buffer 已满。

\section{四种组合状态}

\begin{center}\small
\begin{tabularx}{\linewidth}{L{2.2cm}Y Y}
\toprule
更紧约束 & 可观察证据 & 正确动作/解释\\
\midrule
$rwnd<cwnd$ & advertised window 小、buffer/读取事件 & 受 receiver 限制；等待消费/窗口更新\\
$cwnd<rwnd$ & 拥塞算法状态更小 & 受 path 限制；按 NET07 响应反馈\\
$rwnd=cwnd$ & 数值相等 & 不能仅凭相等判断成因，仍保留两个 Owner\\
任一为 0 & zero receive window 或拥塞状态 & 必须看是哪一状态为 0，恢复机制不同\\
\bottomrule
\end{tabularx}
\end{center}

窗口扩大也不必立刻产生发送：还需要应用有待发数据、连接状态允许、FlightSize 留有空间。窗口缩小不自动撤回已经在途的数据；做题应按题设状态机处理，而不是把新上限追溯应用于已发送事实。

\section{事件账本与调用协议}

\begin{enumerate}
\item 记录 $rwnd,cwnd,FlightSize$ 与待发数据；
\item 标注事件来源：receiver ACK/window、path ACK/loss/ECN/timer 或 application read/write；
\item 只更新该 Owner 的状态，再计算新的 $\min$ 上限；
\item 减去 FlightSize 得 usable；不足一个题设发送单位时停止新发；
\item 若 zero window，进入 probe 分支；若 congestion signal，进入 NET07 算法分支。
\end{enumerate}

\begin{warnbox}{高频 First Divergence}
直接写 $W=rwnd+cwnd$，是把两个上限误作可叠加资源；只算 $\min$ 不减 FlightSize，是把总许可当新增许可；窗口变小就断言丢包，是忽略 receiver feedback；收到 ACK 就只增 \code{cwnd}，是漏掉 ACK 同时可能推进可靠性与携带新 \code{rwnd}，但各状态仍由不同规则更新。
\end{warnbox}

\begin{corebox}{Compression}
谁容量不足？证据来自 receiver 还是 path？哪个状态改变？已有多少在途？合成后还能新发多少？回答完即分别回到 NET06 或 NET07。
\end{corebox}
\end{document}
